\subsection{Naive Approach}\label{sec:naive}
A straightforward approach to address the problem of re-solving recurrent path conditions is to store all previously solved path conditions and their assignments while testing earlier versions and reuse them when testing later versions. Specifically, the naive approach operates as follows:

\vspace{0.5em} \hrule
\begin{algorithmic}[1]\footnotesize
\Procedure{{\Naive}}{$\programs, \budget$}
\Comment{$\programs$: Successive versions of the program}
    \State $\langle \prevdata, \totaltestcases \rangle \gets \langle \emptyset, \emptyset \rangle$
    \For{$\pgm$ \textbf{in} $\programs$}
        \State $\testcases, \newdatas \gets \symb(\pgm, \budget, \prevdata)$   
        \State $\prevdata \gets \prevdata \cup \newdatas$   \Comment{$(\pc, \assignments) \in \newdatas$}
        \State $\totaltestcases \gets \totaltestcases \cup \testcases$
    \EndFor
    \State \Return $\totaltestcases$
    \EndProcedure
\end{algorithmic}
\hrule
\vspace{0.5em}
The $\Naive$ procedure accumulates query data $\data$ generated while testing successive program versions. Recall that each element in $\data$ is a pair consisting of a path condition ($\pc$) and its satisfying assignments ($\assignments$). For example, if $\pc$ is $\myset{(x = 1) \land (y > 0)}$, the SMT solver returns a satisfying assignment $\myset{x \mapsto 1, y \mapsto 1}$, and the corresponding query-data pair is recorded as:
\[
(\myset{(x = 1) \land (y > 0)}, \myset{x \mapsto 1, y \mapsto 1}) \in \data.
\]
If $\pc$ is unsatisfiable, the solver returns an empty assignment set, and the pair $(\pc,\emptyset)$ is added to $\data$. After testing each program version, the generated query data are accumulated into $\prevdata$.

Using the accumulated query data $\prevdata$, the naive approach avoids re-solving path conditions that have already been encountered in previous versions. Specifically, when $\prevdata$ is provided as input to Algorithm~\ref{alg:symbexec}, the $\Solve$ function first checks whether $\prevdata$ contains a pair whose path condition is $\pc \wedge \bc$. If $\prevdata$ contains a pair $(\pc \wedge \bc,\assignments)$, $\Solve$ directly returns the stored result according to the following rule:
\[
\langle \truesat,\assignments\rangle =
\begin{cases}
\langle true,\assignments\rangle
& \text{if } \assignments \neq \emptyset,\\
\langle false,\assignments\rangle
& \text{otherwise}.
\end{cases}
\]
Otherwise, $\Solve$ falls back to the SMT solver and returns the computed result.


\begin{table}[]
    \centering
    \caption{Limitations of naive approach.}
    \label{tab:naive_limitation}
    \resizebox{\columnwidth}{!}{%
    \setlength{\tabcolsep}{4pt}
    \begin{tabular}{l||r|rrr|rr}
    \toprule \midrule
    \multicolumn{1}{c||}{\multirow{2}{*}{\textbf{Program}}} & \multicolumn{1}{c|}{$\prevtwoversion$} & \multicolumn{3}{c|}{$\prevversion$}                                                                                    & \multicolumn{2}{c}{$\newversion$}                                                \\ \cmidrule{2-7} 
    \multicolumn{1}{c||}{}                           & \multicolumn{1}{c|}{\textbf{Store Cost}}   & \multicolumn{1}{c}{\textbf{Match Cost}} & \multicolumn{1}{c}{\textbf{TC Overlap}} & \multicolumn{1}{c|}{\textbf{Store Cost}} & \multicolumn{1}{c}{\textbf{Match Cost}} & \multicolumn{1}{c}{\textbf{TC Overlap}} \\ \midrule \midrule
    diff                                            & 23.6G                                      & \multicolumn{1}{r}{12,977.9}            & \multicolumn{1}{r}{75.6\%}              & 73.9G                                    & \multicolumn{1}{r}{16,450.4}            & 94.5\%                                  \\
    du                                              & 1.2G                                       & \multicolumn{1}{r}{25,249.2}            & \multicolumn{1}{r}{67.1\%}              & 8.7G                                     & \multicolumn{1}{r}{28,817.2}            & 77.6\%                                  \\
    gawk                                            & 4.2G                                       & \multicolumn{1}{r}{9,616.5}             & \multicolumn{1}{r}{88.7\%}              & 6.5G                                     & \multicolumn{1}{r}{13,028.7}            & 94.5\%                                  \\ 
    objdump                                         & 2.4G                                       & \multicolumn{1}{r}{35,087.2}            & \multicolumn{1}{r}{75.6\%}              & 3.8G                                     & \multicolumn{1}{r}{38,026.3}            & 94.5\%                                  \\ \midrule
    average                                         & 7.9G                                       & \multicolumn{1}{r}{20,732.7}            & \multicolumn{1}{r}{76.7\%}              & 23.2G                                    & \multicolumn{1}{r}{24,080.6}            & 90.3\%                                  \\ \midrule \bottomrule
    \end{tabular}%
    }
\end{table}

However, the naive approach suffers from three major limitations: (1) the storage cost of maintaining query data, (2) the cost of matching newly generated path conditions against the stored query data, and (3) the generation of redundant test cases caused by reusing previously stored assignments. Table~\ref{tab:naive_limitation} reports the memory required to store query data, the time spent matching path conditions against the stored query data, and the percentage of generated test cases that are identical to those produced for previous versions.

On average, storing query data from $\prevtwoversion$ requires 7.9 GB of memory. When these data are reused while testing $\prevversion$, path-condition matching consumes 5.8 hours of the 12-hour budget, accounting for 74.6\% of the total SMT-solving time. Furthermore, 76.7\% of the generated test cases are duplicates of those produced for $\prevtwoversion$. These limitations become increasingly severe as testing proceeds across successive releases. After testing $\prevversion$, the storage cost grows to 23.2 GB. In $\newversion$, the matching cost and test case overlap increase to 6.7 hours and 90.3\%, respectively, representing 16.1\% and 17.7\% increases.

Although the naive approach benefits from reusing previously solved path conditions, these benefits are gradually offset by the growing costs of storing query data, matching path conditions, and generating redundant test cases. As shown in Figure~\ref{fig:motivation}, the naive approach initially covers more branches than standard symbolic execution when testing $\prevversion$, where the storage and matching costs remain relatively low. Recall that standard symbolic execution tests each version independently without reusing query data. However, as query data accumulate across releases, the coverage advantage of the naive approach gradually diminishes. In the case of \texttt{diff}, which incurs the largest storage cost, the overheads eventually outweigh the benefits of reuse, causing the naive approach to cover fewer branches than standard symbolic execution.

These observations motivate a new technique specialized for testing successive versions of the same program. The technique should efficiently refine query data collected from prior versions and reuse the refined data when testing subsequent versions. By doing so, it can reduce redundant constraint solving while generating diverse test cases, ultimately improving code coverage and bug detection.